\input includes/\format

\pdfinfo{
   /Title (Gokyo Shumyo)
   /Creator (TeX)
   /Author (Vit Brunner)
   /CreationDate (20261111233500)
}

~
\title
\parindent=0em
\baselineskip=1em

\vskip 5cm
\centerline{\japanesetitle 碁経衆妙}
\bigskip
\centerline{Gokyo Shumyo}
\vfill\break

\input includes/\format-motto

\bigskip
{\header Preface}
\medskip
Gokyo Shumyo is a classical problem collection published by Hayashi Genbi in 1822. Its name could be translated to English as ``Brilliances from go classics''\negthinspace. The collection consists of 520 problems. All positions are shown as {\bf black to play}.
\bigskip

\bigskip

\def\chapterentry#1#2#3{%
  \hbox to \hsize{\strut
    \pdfstartlink attr{/Border [0 0 0]} goto name {problem-#1-1}%
    #1. #2 (#3 problems)\dotfill \chapterpage{#1}%
    \pdfendlink}}
\vbox{
\noindent{\header Sections}
\medskip
\chapterentry{1}{Living}{103}
\chapterentry{2}{Killing}{71}
\chapterentry{3}{Creating a ko}{90}
\chapterentry{4}{Capturing races}{96}
\chapterentry{5}{Oiotoshi}{40}
\chapterentry{6}{Connecting}{74}
\chapterentry{7}{Various techniques}{46}
}

\bigskip

\bigskip
\vbox{
\noindent{\header Third edition note}
\medskip
\noindent Originally published 2006, it's been twenty years. This booklet has been downloaded over a hundred thousand times. I never imagined the reach it would have.

\rightline{\it\me}
\rightline{\it November 2026}
}
\vfill\break

\input includes/\format-gotc
\baselineskip=1.1em
\input generated/\format-gokyoshumyo

\bye
